\section{Evaluation}
\label{sec:evaluation}

\subsection{Exact outcome evaluator}
The primary evaluator uses the hidden site only after an action is submitted. Final simple
regret is
\begin{equation}
 \Rfinal(\theta,\hat x)=f_\theta(x^\star_\theta)-f_\theta(\hat x),
 \qquad x^\star_\theta\in\arg\max_{x\in\mathcal X}f_\theta(x).
 \label{eq:final-regret}
\end{equation}
Experimental opportunity cost charges every occupied unit for the objective it forwent,
\begin{equation}
 \Rcum=\Delta\sum_{r=1}^R\sum_{u\in\mathcal U_r}
 \left[f_\theta(x^\star_\theta)-f_\theta(x_u)\right].
 \label{eq:cumulative-regret}
\end{equation}
Equations~\ref{eq:final-regret} and~\ref{eq:cumulative-regret} are exact up to numerical optimisation of the hidden response surface and support the model and tool comparisons. They are reported separately: an agent may find a
good recommendation through an economically poor campaign.

\subsection{Common-history diagnostic}
Outcome regret alone cannot say whether a failure came from missing evidence or poor use of
available evidence. For the history visible at time $t$, define remaining Bayes risk
\begin{equation}
 b(H_t)=\min_{x\in\mathcal X}\mathbb E[\Rfinal(\theta,x)\mid H_t]
 \label{eq:bayes-risk}
\end{equation}
and recommendation excess risk
\begin{equation}
 g(H_t,\hat x_t)=
 \mathbb E[\Rfinal(\theta,\hat x_t)\mid H_t]-b(H_t)\geq0.
 \label{eq:excess-risk}
\end{equation}
In Equations~\ref{eq:bayes-risk} and~\ref{eq:excess-risk}, large $b(H_t)$ indicates that the visible history still leaves a difficult decision, while large $g(H_t,\hat x_t)$ indicates that the recommendation is worse than the best action supported by that same history. Both condition on identical information and therefore avoid the
Version~1 comparison between a preposterior design score and a realised recommendation.

The diagnostic evaluator uses independent posterior samples for action selection and
evaluation. Crop and economic factors update separately until terminal revenue couples
them. We report diagnostics only where calibration error is smaller than the interpreted
effect; all primary claims remain based on exact hidden-site outcomes.
